\section{C.7. Periods, compact subgroups, and the transcendence question}

We use a Weierstrass equation over $\overline\Q$ and its invariant
differential. Its period lattice is $\Lambda=\Z\omega_1+\Z\omega_2$,
with $\Im(\omega_2/\omega_1)>0$. The quasi-periods are defined by
$\zeta(z+\omega_j)=\zeta(z)+\eta_j$. We extend
$\eta:\Lambda\to\C$ real-linearly to $\C$. The discriminant in
this section is the discriminant of this equation, equivalently of
the lattice in the specified differential. The quantity $N$ is
independent of changing this algebraic equation and differential.

\begin{statement}[C.7.1; Proved: sigma expression and the third-kind period]
Let $P\in E(\overline\Q)$, with $P\ne O$ and $2P\ne O$, and
let $u$ be an elliptic logarithm of $P$. Then
\begin{equation}\label{eq:cp-sigma}
 2\log N(P)=\frac16\log|\Delta|
       +\Re\{\Lambda_1+u(\zeta(u)-\eta(u))\},
 \qquad \Lambda_1=2\log\sigma(u)-u\zeta(u).
\end{equation}
The choices of logarithms change the expression inside the real
part by purely imaginary integral multiples of $2\pi$.
There is an extension $G_P$ of $E$ by $\mathbb G_m$, defined
over $\overline\Q$ and corresponding, under the principal
polarization, to $P$, for which $\Lambda_1$, modulo a logarithm
of an algebraic number, is the negative of a relative period
of $[\Z\to G_P]$. Changing an algebraic differential of the
third kind adds an algebraic multiple of $u$ and algebraic
endpoint terms.
\end{statement}
\begin{proof}
The sigma--theta product, or its divisor and automorphy law
together with its derivative at zero, gives
\begin{equation}\label{eq:cp-klein}
 N(P)=|\Delta|^{1/12}|\sigma(u)|
                   \exp\!\left(-\frac12\Re(u\eta(u))\right).
\end{equation}
To check the constant, for $\Lambda=\Z+\tau\Z$ one uses
$\sigma(z)=e^{\eta_1z^2/2}\vartheta_1(z,\tau)/(2\pi\eta(\tau)^3)$,
$\Delta=(2\pi)^{12}\eta(\tau)^{24}$, and
$\eta(z)=\eta_1z-2\pi i\Im z/T$. Substitution gives exactly
the defining Gaussian in $N$. Homogeneity gives the general
lattice formula. Taking logarithms proves \eqref{eq:cp-sigma}.

We make the period identification explicit. Put
\[
 f_u(z)=\frac{\sigma(z+u)}{\sigma(z)\sigma(u)}e^{-\zeta(u)z}.
\]
Its logarithmic differential is
\[
 \frac{df_u}{f_u}
 =\frac12\frac{\wp'(z)-\wp'(u)}{\wp(z)-\wp(u)}\,dz.
\]
This is the pullback of an algebraic differential with residues
$+1$ at $-P$ and $-1$ at $O$. The apparent $0/0$ at $z=u$
has a removable singularity when $P$ is not $2$-torsion.
The rational cocycle $f_u(z+z')/(f_u(z)f_u(z'))$ defines
the algebraic extension; in its rational trivialization the
exponential map has the form
\[
 \exp_{G_P}(z,v)=(\exp_E z,e^v f_u(z)).
\]
Its periods include
\begin{equation}\label{eq:cp-periods}
 (0,2\pi i),\qquad
 (\omega_j,\zeta(u)\omega_j-\eta_j u),\quad j=1,2.
\end{equation}
These signs follow directly from the sigma automorphy law.
The identity
$\sigma(2u)=-\wp'(u)\sigma(u)^4$ implies
\[
 \log f_u(u)=\Lambda_1+\log(-\wp'(u)).
\]
For an algebraic lift $R=(P,r)$, $r\in\overline\Q^{\times}$,
the second coordinate of $\log_{G_P}R$ is
\begin{equation}\label{eq:cp-v}
 v=\log r-\Lambda_1-\log(-\wp'(u)).
\end{equation}
It is a period of the relative homology defining the indicated
$1$-motive. This also avoids treating an integral with an endpoint
at a residue pole as an unregularized ordinary integral.
The description agrees with the generalized logarithm and
Serre function in \citet[Definitions 2.3--2.4 and
Section 6]{bertolin2022}; the signs in
\eqref{eq:cp-periods} are for the exponential rather than the
quasi-periods of $f_u$ itself.
\end{proof}

\begin{statement}[C.7.2; Proved: the compact-subgroup criterion and the real formula]
For $P$ as in Statement C.7.1, $N(P)$ is algebraic if and only
if the fibre of $G_P$ above $P$ contains an algebraic point of
the maximal compact subgroup $G_P(\C)^c$.
For the family in the prompt, if $t>1$ is real, choose
$\tau=iT$, $T>1$, and the zero $\mu=1/2-iy$, $0<y<T/2$.
With $x=2y/T$, one has
\[
 N(P)^4=3^{3/4}(t-1)^{-1/4}
       \left|\frac{\vartheta_1(x,i/T)}{\eta(i/T)}\right|.
\]
\end{statement}
\begin{proof}
On the analytic universal cover of $G_P$ define the real
homomorphism
\[
 \chi(z,v)=\Re\{v-\zeta(u)z+u\eta(z)\}.
\]
It vanishes on the periods \eqref{eq:cp-periods}, hence descends
to $G_P(\C)$. Its restriction to the torus is $\log|\cdot|$.
The kernel is compact: in the universal cover it is a real
three-dimensional vector space modulo the rank-three period
lattice. Every compact subgroup has zero image in $\R$, so
this kernel is exactly $G_P(\C)^c$.
The identity
$\Re(u\eta(z))=\Re(z\eta(u))$ follows from the Legendre relation;
equivalently $\eta(z)=sz+\pi\bar z/\operatorname{area}(\Lambda)$
for a complex constant $s$. Substituting
\eqref{eq:cp-v} at $z=u$ gives
\begin{equation}\label{eq:cp-chi}
 \chi(R)=\log|r|-\log|\wp'(u)|
            -2\log N(P)+\frac16\log|\Delta|.
\end{equation}
If $r$ is algebraic and $\chi(R)=0$, this equation makes
$N(P)^2$ algebraic. Conversely, if $N(P)$ is algebraic,
choose the positive real algebraic number
$r=|\wp'(u)|\,|\Delta|^{-1/6}N(P)^2$.
Then \eqref{eq:cp-chi} vanishes. This proves both directions
with the chosen rational trivialization of the fibre.
The compact interpretation is the one discussed by
\citet[slides 5--8]{bertrand2017}; here the calculation fixes
its normalization on the diagonal $P=P$.

For the real formula, the rectangular elliptic curve has its
two zeros of $\wp$ on $1/2+i\R$; along that edge $\wp$ runs
between the positive root $e_1$ and the negative root $e_2$.
Thus the indicated representative exists. Translate $2\mu$
by $-1$ and apply the modular substitution $\tau\mapsto-1/\tau$.
The resulting elliptic argument is $-2y/T$, which is real.
The Gaussian is therefore one. Evenness of $N$ and the
duplication identity prove the displayed formula. This removes
the Gaussian from this particular expression but makes no
algebraicity assertion about its theta value.
\end{proof}

\subsection{The conditional theorem and its exact conjectural input}

We use Grothendieck's period conjecture in the form stated as
Conjecture 1.1 of \citet{bertolin2022}: for a motive $M$ over
$\overline\Q$, the field generated by its periods over
$\overline\Q$ has transcendence degree $\dim G_{\mathrm{mot}}(M)$.
Only $1$-motives occur below. We also use their Construction 3.4
and Lemma 3.5: if $B$ is the smallest abelian subvariety
generated by the point and extension data, and $Z(1)$ is the
central torus constructed from the Weil-pairing bracket and
the remaining algebraic lift, then
\begin{equation}\label{eq:cp-MTdim}
 \dim G_{\mathrm{mot}}(M)
 =\dim G_{\mathrm{mot}}(A)+2\dim B+\dim Z(1).
\end{equation}
The bracket is described explicitly in their Section 6.1.
Their Table 2 corrects the table in the earlier paper; the
dependent nondeficient case has unipotent dimension three.
We now verify which case occurs here.

\begin{statement}[C.7.3; Proved: the motivic groups needed on the diagonal]
Let $P$ be nontorsion and let $R$ be any algebraic lift above
$P$ in $G_P$. Put $M_R=[\Z\to G_P]$, $1\mapsto R$.
Its connected motivic group is, after a choice of Levi subgroup,
the semidirect product of $G_{\mathrm{mot}}(E)$ and the
three-dimensional Heisenberg group attached to $H_1(E)$ and
its principal polarization. Its dimension is $5$ for CM and
$7$ without CM.

For a non-real embedding put
$\mathcal M=M_R\oplus\overline{M_R}$.
If $E$ and $\bar E$ are not isogenous, its dimension is $13$.
If they are isogenous, choose an isogeny
$\phi:\bar E\to E$ of degree $d$ and write
$Q=\phi(\bar P)$. Let $r$ be the dimension over
$F=\End(E)\otimes\Q$ of the span of $P,Q$, modulo torsion.
Then
\[
\begin{array}{c|c|c|c}
\text{case}&\dim B&\dim Z&\dim G_{\mathrm{mot}}(\mathcal M)\\ \hline
r=2,\ E\text{ CM}&2&2&8\\
r=2,\ E\text{ not CM}&2&2&10\\
r=1,\ E\text{ CM}&1&1+e&5+e\\
r=1,\ E\text{ not CM}&1&1+e&7+e.
\end{array}
\]
Here $e$ is either zero or one: it is zero precisely when the
algebraic lift in the one-dimensional quotient of the two
central tori by the bracket torus is torsion. This criterion is
independent of the choices made modulo isogeny. In the real
case we use the single motive $M_R$; no extra quotient torus
is necessary.
\end{statement}
\begin{proof}
All curves, polarizations, points, and lifts are algebraic over
$\overline\Q$, as required by Construction 3.4 and Lemma 3.5.
For $M_R$ the smallest abelian subvariety containing $(P,P)$ is
the diagonal $B=E$, since $P$ is nontorsion. If $\psi$ is the
polarization pairing, the bracket on this diagonal is
\[
 [(x,x),(y,y)]=\psi(x,y)-\psi(y,x)=2\psi(x,y).
\]
It is nonzero and spans $\Q(1)$. Hence $Z'=Z=\mathbb G_m$;
there is no deficient or Ribet exception, including for CM.
The pure group is $\mathrm{GL}_2$ without CM and
$\operatorname{Res}_{F/\Q}\mathbb G_m$ for the imaginary
quadratic CM field. Their dimensions are four and two.
This proves the single-motive assertion, including its Lie
bracket and Levi action: the action on the centre is the
polarization multiplier.

For the direct sum only the two diagonal central coordinates
occur. If the two elliptic curves are not isogenous, they
cannot be CM conjugates. The pure group is
$\{(g_1,g_2)\in\mathrm{GL}_2^2:\det g_1=\det g_2\}$,
of dimension seven. The point data generate $E\times\bar E$,
and the two brackets are independent, giving $7+4+2=13$.

If the curves are isogenous, transport the second polarized
factor through $\phi$. Its polarization pairing becomes
$\psi/d$. The smallest abelian subvariety has dimension $r$.
For $r=2$, independent variation in its two factors makes the
two diagonal brackets span both central coordinates; thus
$Z'=Z=\mathbb G_m^2$. The unipotent group is the product
of two Heisenberg groups, with the same pure group acting
on both standard factors.
For $r=1$, after multiplying the points by a common integer
we may write $Q=aP$, $a\in F^\times$. The bracket image is
the one-dimensional line with proportional central coordinates
\begin{equation}\label{eq:cp-positive-centre}
 (1,c),\qquad c=\frac{a a^\dagger}{d}>0,
\end{equation}
where $\dagger$ is the Rosati involution. On an elliptic curve
$aa^\dagger$ is a positive rational number: a square in the
non-CM case and the positive field norm in the CM case.
Construction 3.4(3) adds a further central direction exactly
when the image of the algebraic lift in the quotient torus
is nontorsion. This is the integer $e$ above.
The unipotent Lie algebra is the Heisenberg algebra, with one
additional central Tate line when $e=1$. Formula
\eqref{eq:cp-MTdim} proves every entry of the table. This
description computes the group even when the lift contributes
a Kummer factor; specifying only the span of $P,Q$ would
not determine that last factor.
\end{proof}

\begin{statement}[C.7.4; Proved conditionally: transcendence of the local norm]
Assume Grothendieck's period conjecture for $1$-motives over
$\overline\Q$. For every algebraic parameter $t\ne0,1$ at
which $P_t$ is nontorsion, $N(P_t)$ is transcendental.
More generally the conclusion holds for every nontorsion
algebraic point on an elliptic curve over $\overline\Q$.
The weaker named Relative Semi-abelian Conjecture 4.7 of
\citet{bertolin2022}, restricted to the groups and algebraic
lifts used in this proof, already suffices.
\end{statement}
\begin{proof}
Suppose $N(P)$ is algebraic. Statement C.7.2 supplies an
algebraic lift $R=(P,r)$ in $G_P(\C)^c$, and its conjugate
supplies a lift $\bar R$ in $G_{\bar P}(\C)^c$.
Let $v$ denote the logarithmic coordinate in
\eqref{eq:cp-v}. Put
\[
 K_0=\overline\Q\bigl(\text{entries of the full de Rham--Betti
                    period matrix of }G_P\times G_{\bar P}\bigr).
\]
Here ``full'' includes the integrals of differentials of the second
kind, hence the quasi-periods $\eta_j$, as well as the periods
$\omega_j$, the third-kind periods $\kappa_j$, and $2\pi i$.
It is the matrix $\Omega_G$ in their Definition 2.4, not merely
the lattice in the Lie algebra of the semi-abelian variety.
This field contains $u,\zeta(u)$ and their complex conjugates.
Indeed it contains the pure periods, $2\pi i$, and
$\kappa_j=\zeta(u)\omega_j-\eta_j u$; the determinant of
the two equations for $u,\zeta(u)$ is, up to sign,
$\omega_1\eta_2-\omega_2\eta_1= -2\pi i$.
This is also the observation of \citet[Remark 4.13]{bertolin2022}.
The real-linear expression for $\eta(u)$ is
\begin{equation}\label{eq:cp-real-linear}
 \eta(u)=
 \frac{(u\bar\omega_2-\bar u\omega_2)\eta_1
       +(\bar u\omega_1-u\bar\omega_1)\eta_2}
      {\omega_1\bar\omega_2-\bar\omega_1\omega_2}.
\end{equation}
The denominator is nonzero. Thus $u\eta(u)$ and its conjugate
belong to $K_0$. The compact equation $\chi(R)=0$ becomes
the algebraic relation over $K_0$
\begin{equation}\label{eq:cp-relation}
 v+\bar v
 =u\zeta(u)+\bar u\overline{\zeta(u)}
                -u\eta(u)-\bar u\overline{\eta(u)}.
\end{equation}
Complex conjugates here are coordinates of the conjugate
$1$-motive. They are not being treated as rational functions
of un-conjugated periods without justification.

Consider the period torsor of $\mathcal M$ over the period
torsor of $G=G_P\times G_{\bar P}$. The kernel of the
corresponding motivic groups is its weight $-2$ vector group
$\operatorname{Lie}Z$, of dimension $\dim Z$. There is no
weight $-1$ kernel: the point coordinates are exactly the
extension coordinates under the polarization, so $B_Q=0$.
More explicitly, projection of the smallest abelian variety $B$
to the variety generated by the extension coordinates is an
isogeny, since $B$ is the graph of the same point data in the
two factors. It therefore induces an isomorphism on the
weight $-1$ graded Lie algebras. The motive $[0\to G]$ has
only weights $-1$ and $-2$, and its unipotent Lie algebra has
no weight $-2$ graded piece. Thus the entire kernel is exactly
$W_{-2}\operatorname{Lie}G_{\rm mot}(\mathcal M)$; no additional
abelian direction or unaccounted logarithm remains in the fibre.
After evaluating the base periods, the fibre is an affine
torsor over $K_0$ under that vector group. Translation in
its two central coordinates changes $(v,\bar v)$ by
$2\pi i(c_1,c_2)$, with the signs fixed by using a positive
Tate loop for each motive. The linear functional
$c_1+c_2$ is nonzero on this central group. It is nonzero
on the whole plane when $\dim B=2$, and when $\dim B=1$
it is already nonzero on the bracket line
\eqref{eq:cp-positive-centre}, because $1+c>0$.

Consequently \eqref{eq:cp-relation} cuts a proper affine
hyperplane in the period-torsor fibre. It forces the relative
transcendence degree of the actual periods over $K_0$ to be
at most $\dim Z-1$. The pure-base period torsor has dimension
$\dim G_{\mathrm{mot}}(G)$, so the total transcendence degree
is at most $\dim G_{\mathrm{mot}}(\mathcal M)-1$.
This contradicts the asserted period conjecture. For real
data the single motive suffices: conjugation acts on its
Betti realization over $\Q$, and the same equation has a
nonzero coefficient on its one central direction. One can
also double the point and lift into the identity component,
choose $u,\omega_1$ real, and read the relation directly as
\[
 \omega_1 v-\omega_1u\zeta(u)+\eta_1u^2=0.
\]
Doubling preserves the equivalence of algebraicity of $N$
by the division formula and its algebraic nonzero factors.

For the weaker conjectural input, Conjecture 4.7 asserts
relative transcendence degree at least $2\dim B_Q+\dim Z$.
All $R_i$ used here are algebraic, $B_Q=0$, and the
generalized logarithms add exactly the central coordinates
to $K_0$. Thus it contradicts the same hyperplane relation
directly. Proposition 4.12 of that paper derives Conjecture
4.7 from the generalized period conjecture; in our case the
field of definition is algebraic, so the ordinary period
conjecture is enough. The weaker Conjecture 4.9, formulated
using arbitrary algebraic translates, gives a zero lower
bound at an algebraic point and does not by itself yield
this conclusion. Linear relations among $1$-periods are
also insufficient: after clearing the denominator in
\eqref{eq:cp-real-linear}, the prohibited relation is
polynomial, not a linear relation with algebraic coefficients.
\end{proof}

\begin{statement}[C.7.5; Proved computationally: the specified integer-relation tests]
At all eleven rational CM parameters, an independent theta
calculation and the global height calculation agree to more
than $245$ decimal places. At $200$ decimal digits, PSLQ
with tolerance $10^{-185}$, maximum coefficient $10^8$ and
$10000$ steps returns no relation in each prescribed list
consisting of $\log N$, logarithms of the bad primes together
with $\log2,\log3$, $\pi$, and independent real period
components. At $j=287496$ the known identity
$\omega_1=2|\Im\omega_2|$ is removed first; repeated prime
logarithms and zero period components are likewise omitted.
An additional direct algebraicity search tests both $N$ and $N^{12}$
at all eleven points, using the vectors $(1,x,\ldots,x^d)$ for
$d=4,8,12$, working precision $235$ decimal digits, tolerance
$10^{-210}$, maximum coefficient $10^{12}$ and $20000$ steps.
All $66$ calls return no relation.
\end{statement}
\begin{proof}[Computational verification]
The complete inputs, $255$-digit values, comparison errors,
PSLQ labels and returns are in
\path{code/C_cm_250.gp}, \path{code/C_cm_250.out},
\path{code/C_cm_theta_pslq.py}, and its output. The theta
calculation uses the exact CM period ratios and solves the
Fourier equation for $\wp(1/2-iy)=0$ independently of PARI's
height implementation. It verifies the singular modulus
as well. The additional polynomial search is reproduced by
\path{code/C7_norm_polynomial_search.py}; its output and JSON
record include every parameter, source hash and returned vector.
These tests prove no transcendence assertion. A null PSLQ return
is not an exhaustive certificate excluding all polynomials within
the stated coefficient and degree limits.
\end{proof}

\subsection{The unconditional problem remains under investigation}

Statements C.7.1--C.7.5 do not decide the unconditional
arithmetic nature of $N$ at a nontorsion point. An explicit
small instance is
\[
 E:y^2=x^3-140x+784,\qquad P=(0,28),\qquad j=-3375,
 \qquad N(P)=2e^{-\widehat h(P)}.
\]
Here $P$ is a primitive generator modulo the rational
$2$-torsion and
\[
 \widehat h(P)=
 0.46736567718147128000752020126284678363562034370331\ldots.
\]
The exact unresolved assertion is that $2e^{-\widehat h(P)}$
is transcendental. The corresponding example of conductor
$361$, $y^2+y=x^3-38x+90$, $P=(0,9)$, has
$N=e^{-\widehat h(P)}$ and no finite local contribution.
The relative-period proof isolates its obstruction: an
unconditional theorem excluding
\eqref{eq:cp-relation} for the positive diagonal central
direction. A theorem about linear $1$-period relations
does not supply that exclusion. Further attempts and the
CM theta comparison are recorded below as they are tested.
